% ============================================================
% From Local Curvature to the Weil Functional:
% An Explicit Construction
%
% Author:   Ulrich Tehrani
% Date:     May 2026
% Version:  36
% Zenodo:   https://doi.org/10.5281/zenodo.19106992
% License:  CC BY 4.0
% ============================================================

\documentclass[11pt]{article}
\usepackage{cmap}
\usepackage[T1]{fontenc}
\usepackage[utf8]{inputenc}
\usepackage[a4paper,margin=1in]{geometry}
\usepackage{amsmath,amssymb,amsthm,mathtools}
\usepackage{xcolor}
\usepackage{hyperref}
\usepackage{enumitem}
\usepackage{setspace}
\usepackage{booktabs}
\usepackage{array}
\usepackage{graphicx}

\hypersetup{colorlinks=true,
  linkcolor=blue!50!black,
  urlcolor=blue!50!black,
  citecolor=blue!50!black}
\setlength{\parskip}{0.5em}
\setlength{\parindent}{0pt}
\onehalfspacing
\emergencystretch=2em

% --------------------------------------------------------
% Theorem environments (numbered within section)
% --------------------------------------------------------
\newtheorem{theorem}{Theorem}[section]
\newtheorem{proposition}[theorem]{Proposition}
\newtheorem{lemma}[theorem]{Lemma}
\newtheorem{corollary}[theorem]{Corollary}
\newtheorem{definition}[theorem]{Definition}
\newtheorem{assumption}[theorem]{Assumption}
\newtheorem{observation}[theorem]{Observation}
\newtheorem{conjecture}[theorem]{Conjecture}
\newtheorem*{remark}{Remark}
\newtheorem{openproblem}[theorem]{Open Problem}

% --------------------------------------------------------
% Series-wide macros
% --------------------------------------------------------
\newcommand{\Hlocal}{H_{\mathrm{local}}}
\newcommand{\Hren}{H_{\mathrm{local}}^{\mathrm{ren}}}
\newcommand{\Sad}{S_{\mathrm{ad}}}
\newcommand{\cren}{c_p^{\mathrm{ren}}}
\newcommand{\Zren}{Z_{\mathrm{ren}}}

% --------------------------------------------------------
% Title
% --------------------------------------------------------
\title{\vspace{-1.5em}\bfseries
From Local Curvature to the Weil Functional:\\[0.4em]
\large An Explicit Construction}
\author{Ulrich Tehrani}
\date{Research Note \textperiodcentered\ May 2026
      \textperiodcentered\ v36}

% ============================================================
\begin{document}
\maketitle
% ============================================================

\begin{abstract}
Building on the curvature decomposition established
in~\cite{Paper1}, we construct an explicit family of admissible
test functions $g^*_{\sigma,\varepsilon}\in\Sad$ for the Weil
explicit formula and identify a formal diagonal localisation
pattern that targets the local curvature
$\Hlocal(\sigma,\kappa)$ on the prime side.
A renormalised family of prime weights
$\cren=f_p^{1/2}>0$ is introduced, and the associated diagonal
energy $D=\sum_p (\cren)^2\approx 9.470$ is shown to converge.
The renormalised local curvature is shown to be of lower order
than $(\log\kappa)^2$ and to exhibit a sawtooth structure related
to classical Mertens-type prime-sum asymptotics.
No RH, GUE, or Hilbert--P\'olya input is used.
No theorem or result in this paper is conditional on
positive answers to these open problems: the paper isolates three open
questions for the Weil-normalised embedding.

\textbf{What is proved.}
The antisymmetrised Gaussian family $g^*_{\sigma,\varepsilon}$
belongs to the admissible class $\Sad$ used in
Lagarias~\cite{Lagarias} (Lemma~\ref{lem:admissible}).
The renormalised prime weight
$\cren := \bigl(V_p(\tfrac12)-\tfrac{4(\log p)^2}{p}\bigr)^{1/2}$
is strictly positive (Lemma~\ref{lem:fp-pos}), and the diagonal
energy series $D=\sum_p (\cren)^2$ converges absolutely
(Lemma~\ref{lem:D-conv}).
The normalised Gaussian correlation $K_\varepsilon$
localises the diagonal first-prime contribution exactly
to $V_p(\sigma)$ by an algebraic identity
(Observation~\ref{obs:diagonal}).

\textbf{What is numerical.}
The diagonal energy computation gives
$\sum_{p<10^7} f_p = 9.4702$, with tail
$\sum_{p\ge 10^7}f_p < 2.7\times 10^{-4}$, hence
$D\approx 9.470$ to three decimal places.
Separately, at zero-sum reference parameters
$(\varepsilon,N)=(0.05,100)$, direct grid evaluation of the
prime-only renormalised quantities at
$\kappa\in\{10,20,29,50,100,200,500,1000\}$ shows that
$\Zren\ge\Hlocal^{\mathrm{pr}}$ holds at every tested point
except $\kappa=20$,
where $\Zren-\Hlocal^{\mathrm{pr}}=-0.988$, with the ratio
$\Zren/\Hlocal^{\mathrm{pr}}$
equal to about $1.5$ at $\kappa=29$ and about $10.1$ at $\kappa=1000$
(prime-only comparison; the archimedean term
$\psi_1(\tfrac12)\approx 2.150$ enters separately).
No claim is made about $\kappa$ outside the eight tested points.

\textbf{What remains open.}
Three questions separate the present construction from a full
prime-side derivation of Weil positivity for this family
(\S\,\ref{sec:open}).
Problem~\ref{op:weil} (\emph{Weil-normalised embedding})
asks whether the diagonal localisation of
Observation~\ref{obs:diagonal} can be implemented inside the
Lagarias--Weil functional, including the correct Gaussian
rescaling, the archimedean contribution, and the zero-side
contribution in its unconditional form.
Problem~\ref{op:spectral-ineq} (\emph{Spectral inequality})
asks whether, for every fixed cutoff~$\kappa$, there exists
$\varepsilon_0(\kappa)>0$ such that
$\mathcal{Z}_W(g^*_{1/2,\varepsilon})
\ge\Hlocal(\tfrac12,\kappa)$ holds unconditionally for all
$0<\varepsilon<\varepsilon_0(\kappa)$.
Problem~\ref{op:offdiag} (\emph{Off-diagonal control}) asks for
an unconditional bound on the cross-terms for the proxy family.

\textbf{Structure.}
The formal core in \S\S\,\ref{sec:construct}--\ref{sec:scale} is
non-circular: the analytic statements are labelled as proved,
the finite computations are labelled as numerical, and the
Weil-normalised embedding questions are isolated as open problems
in~\S\,\ref{sec:open}.
No use of the Riemann Hypothesis, the GUE conjecture, or the
Hilbert--P\'olya postulate is made anywhere in
\S\S\,\ref{sec:construct}--\ref{sec:scale}.
\end{abstract}

\vspace{2em}

% --------------------------------------------------------
% Notation and Conventions
% --------------------------------------------------------
\section*{Notation and Conventions}
\label{sec:notation}

\emph{Critical line.}
Throughout, the critical line means $\Re(s)=\tfrac12$.
The construction parameter $\sigma>0$ is evaluated at
$\sigma=\tfrac12$ in the renormalised analysis unless stated
otherwise.

\medskip
\emph{Global parameters.}
The reference values are $\kappa\ge 2$ (prime cutoff,
varied across the comparison grid),
$\varepsilon=0.05$ (Gaussian smoothing parameter), and
$N=100$ (number of Riemann ordinates $\gamma_k$ used in the
zero-sum evaluations of~\S\,\ref{sec:numerical}).
Values are fixed within each statement and explicitly varied
otherwise.

\medskip
\emph{Phasor convention.}
Throughout this paper the prime phasor is
$\sin(\gamma_k \log p)$ evaluated at $\sigma = \tfrac12$;
no additional $\sigma$-factor enters the Fourier side.
The Fourier transform $\widehat g$ of a test function $g$
uses the convention of Lagarias~\cite{Lagarias};
the admissible test function $g^*_{\sigma,\varepsilon}$ is
defined in log-prime position space
in~\S\,\ref{sec:construct}; its Fourier transform is derived
there.

\medskip
\emph{Main objects.}
The local curvature, imported from~\cite{Paper1}, is
\[
\Hlocal(\sigma,\kappa) \;=\; \psi_1(\sigma)
\;+\; \sum_{p\le\kappa} V_p(\sigma),
\qquad
V_p(\sigma) \;=\; 4(\log p)^2\,
\frac{p^{-2\sigma}}{(1-p^{-2\sigma})^2}.
\]
The Gaussian bump is $\varphi_\varepsilon(x)=
(\varepsilon\sqrt{2\pi})^{-1} e^{-x^2/(2\varepsilon^2)}$.
The raw weight at level $\sigma$ is
$c_{p,\sigma} = V_p(\sigma)^{1/2}\cdot p^{1/4}/(\log p)^{1/2}$,
and the antisymmetrised admissible test function
$g^*_{\sigma,\varepsilon}\in\Sad$ is constructed
in~\S\,\ref{sec:construct}.
The Weil quadratic functional is denoted $W$;
its zero-side contribution in the Lagarias
normalisation~\cite{Lagarias} is denoted $\mathcal{Z}_W(g)$.
The \emph{numerical ordinate proxy} is the finite positive sum
$Z_{\mathrm{ord}}(g;N) := \sum_{k=1}^{N}
|\widehat g(\gamma_k)|^2$
over the first $N$ verified critical-line ordinates $\gamma_k$.
The unconditional Lagarias--Weil zero-side contribution
$\mathcal{Z}_W(g)$ is distinct from $Z_{\mathrm{ord}}$;
the two objects are not identified without an explicit
embedding statement (Problem~\ref{op:weil}).
At $\sigma=\tfrac12$ the renormalised prime weight is
$\cren = f_p^{1/2}$, where
$f_p = V_p(\tfrac12)-4(\log p)^2/p$, and the diagonal
energy is $D=\sum_p (\cren)^2$.
The renormalised local curvature is
$\Hren(\tfrac12,\kappa)=\Hlocal(\tfrac12,\kappa)-2(\log\kappa)^2$.

\medskip
\emph{Analytical foundations.}
The construction uses the Weil explicit formula~\cite{Weil} in the
Lagarias admissible formulation~\cite{Lagarias} and the
prime number theorem for the asymptotic scale of~$\Hlocal$.

\medskip
\emph{Numerical computations.}
All numerical results were computed with \texttt{mpmath}
at 50~decimal digits of working precision;
Riemann zeros $\gamma_k$ were evaluated via
\texttt{mpmath.zetazero}.
Verification code is available at
\url{https://github.com/utehrani/analysislab-nt}.

\medskip
The diagonal energy $D = \sum_p f_p \approx 9.470$ defined here is
specific to the renormalised prime-side construction of this paper.
It is the convergent prime-power-tail constant in~\eqref{eq:D} and
should not be read as a general trace-formula diagonal term.
The renormalised weight $\cren=f_p^{1/2}$ used in the Weil
functional differs in normalisation from the unrenormalised
weight $c_{p,\sigma}$ used in the construction of
$g^*_{\sigma,\varepsilon}$.

\newpage

% --------------------------------------------------------
% §1 Introduction
% --------------------------------------------------------
\section{Introduction}
\label{sec:intro}

\subsection{Background and Motivation}

This paper is the second in a series developing an
operator-theoretic framework for studying the distribution of
the non-trivial zeros of the Riemann zeta function $\zeta(s)$.
The central organising principle of the series is that prime
numbers mediate coupling between zero ordinates, and that
explicit finite-cutoff constructions on the prime side carry
arithmetic information accessible by classical means.

\textbf{Paper~1}~\cite{Paper1} introduced the curvature field
\[
H_\xi(\sigma,t) \;:=\; \frac{d^2}{d\sigma^2}\,
\log|\xi(\sigma+it)|^2
\;=\; \Hlocal(\sigma,\kappa) \;+\; H_{\mathrm{dual}}(\sigma,t,\kappa),
\]
where $H_{\mathrm{dual}}=H_\xi-\Hlocal$ is the bookkeeping
remainder defined as the finite-cutoff
difference (cf.~\cite{Paper1}), and where the truncated
local curvature
\[
\Hlocal(\sigma,\kappa) \;=\; \psi_1(\sigma)
\;+\; \sum_{p\le\kappa} V_p(\sigma)
\]
contains the archimedean term and the finite-place
contributions.
The prime-side pieces satisfy $V_p(\sigma)\ge 0$, and on the
critical line one has $\Hlocal(\tfrac12,\kappa)\asymp
(\log\kappa)^2$.

\textbf{The present paper} addresses a concrete part of the
central question raised in the Open Problems section
of~\cite[\S\,8]{Paper1}: whether the
local positivity and critical divergence of $\Hlocal$ can be
organised into an explicit admissible test-function construction
whose prime-side term in the Weil functional matches the local
curvature.
The present paper addresses the first, purely algebraic part
of this question.
We construct an explicit family
$g^*_{\sigma,\varepsilon}\in\Sad$ and show that its
normalised Gaussian diagonal contribution reproduces each local
curvature term $V_p(\sigma)$ at fixed cutoff
(Observation~\ref{obs:diagonal}).
The implementation of this localisation inside the fully
normalised Lagarias--Weil functional --- including Gaussian
rescaling, the archimedean term, and the unconditional
zero-side contribution --- is left as
Problem~\ref{op:weil}.

The construction $g^*_{\sigma,\varepsilon}\in\Sad$ developed here
should be distinguished from the established operator-theoretic
programmes addressing the Riemann Hypothesis.
It is not a Hilbert--P\'olya spectral realisation in the sense of
Connes~\cite{Connes} or Berry--Keating~\cite{BerryKeating}, nor a
de~Branges structure-function construction~\cite{deBranges}; it
is an explicit antisymmetrised Gaussian smearing of prime delta
masses with arithmetic weights, closest in spirit to Hilbert-space
reformulations of the explicit formula such as those of
Burnol~\cite{Burnol}.
In the formal curvature dictionary of~\cite{Paper1}, the local
curvature $\Hlocal(\sigma,\kappa)$ is the finite-cutoff local
object that the present construction is designed to realise on
the prime side of a Weil-functional framework.
This paper proves the normalised diagonal localisation of the
finite-prime contribution (Observation~\ref{obs:diagonal}); the fully
normalised Lagarias--Weil embedding, including the archimedean
and zero-side terms, is left as Problem~\ref{op:weil}.
The curvature label reflects the
second-derivative origin of $V_p(\sigma)$ in $\log|\xi|^2$
identified in~\cite{Paper1} and is not used in any further
geometric sense.
The terms ``prime--zero coupling'' and ``curvature field'' used
informally in the series are translations of the explicit-formula
duality and of $\partial_\sigma^2\log|\xi|^2$ respectively, not
new mathematical structures.

The construction is set within the Lagarias formulation of Weil
positivity~\cite{Lagarias}, which characterises the Riemann
Hypothesis as positivity of the Weil functional $W(f*\check f)$
on a suitable class of test functions, building on the Weil
explicit formula~\cite{Weil} and the Li positivity
criterion~\cite{Li}.
The renormalised diagonal scale $D=\sum_p(\cren)^2$ encountered
in~\S\,\ref{sec:renorm} is closely related, via the prime number
theorem and Mertens' classical estimates~\cite{Mertens}, to the
average behaviour of the local curvature; this connection is
made precise in~\S\,\ref{sec:scale}.

\subsection{Main Results}

The paper contains three main components.

The first is the explicit construction of the test family
$g^*_{\sigma,\varepsilon}$ and the proof that it lies in the
admissible class $\Sad$ used by Lagarias~\cite{Lagarias}.
The construction proceeds in two steps: a Gaussian smearing of
delta masses at $\log p$ with weights $c_{p,\sigma}$, followed
by antisymmetrisation enforcing the vanishing condition at the
origin (\S\,\ref{sec:construct}).
The admissibility lemma (Lemma~\ref{lem:admissible}) is a
direct verification.

The second is the formal diagonal localisation
(Observation~\ref{obs:diagonal}): the normalised Gaussian
correlation $K_\varepsilon$ targets the local curvature
$V_p(\sigma)$ at each prime by an algebraic identity in the
weights $c_{p,\sigma}$, while all off-diagonal and higher
prime-power contributions are exponentially suppressed as
$\varepsilon\to 0$ at fixed $\kappa$.
The precise embedding of this localisation into the
Lagarias--Weil functional --- including the correct Gaussian
rescaling, the zero-side contribution, and the archimedean
term --- is formulated as an open problem
(Problem~\ref{op:weil}).

The third is the renormalised analysis at the critical line
(\S\S\,\ref{sec:renorm}--\ref{sec:scale}).
At $\sigma=\tfrac12$ the auxiliary quantity
$f_p=V_p(\tfrac12)-4(\log p)^2/p$ is shown to satisfy the
closed form
\[
f_p \;=\; \frac{4(\log p)^2(2p-1)}{p(p-1)^2} \;>\; 0
\qquad (p\ge 2)
\]
(Lemma~\ref{lem:fp-pos}), so that $\cren=f_p^{1/2}$ is real
and strictly positive.
The diagonal energy $D=\sum_p(\cren)^2$ converges absolutely by
comparison with $\sum_n (\log n)^2/n^2$ (Lemma~\ref{lem:D-conv})
and evaluates to $D\approx 9.470$.
The renormalised local curvature
$\Hren(\tfrac12,\kappa)=\Hlocal(\tfrac12,\kappa)-2(\log\kappa)^2$
is of lower order than $(\log\kappa)^2$ and exhibits a sawtooth
structure: a positive jump $V_p(\tfrac12)$ at each new prime
$p\le\kappa$, decreasing continuously between primes as
$2(\log\kappa)^2$ grows.
By contrast, the prime-power-tail energy
$H_{\mathrm{tail}}(\kappa):=\sum_{p\le\kappa}f_p$ jumps
by~$f_p$ at each new prime and converges monotonically to~$D$.
The classical estimate
$\sum_{p\le\kappa}(\log p)/p=\log\kappa+O(1)$ controls
the average drift, with the Mertens constant $\gamma_M$
appearing in the related harmonic prime-sum estimate
$\sum_{p\le\kappa}1/p = \log\log\kappa + \gamma_M + o(1)$.

The numerical comparison of~\S\,\ref{sec:numerical} finds
that $\Zren\ge\Hlocal^{\mathrm{pr}}$ holds at every point of the eight-point
test grid except $\kappa=20$, with the ratio equal to about $1.5$
at $\kappa=29$ and about $10.1$ at $\kappa=1000$.
Three open problems remain:
the Weil-normalised embedding (Problem~\ref{op:weil}),
the unconditional spectral inequality
$\mathcal{Z}_W(g^*_{1/2,\varepsilon})\ge\Hlocal(\tfrac12,\kappa)$
(Problem~\ref{op:spectral-ineq}), and
the off-diagonal control of the cross-terms
(Problem~\ref{op:offdiag}).

\subsection{Reader Contract}

This paper \emph{proves}, by direct construction and classical
identities, the admissibility of the test family
(Lemma~\ref{lem:admissible}), the strict positivity
of the renormalised weight $\cren=f_p^{1/2}>0$
(Lemma~\ref{lem:fp-pos}), and the absolute convergence of the
diagonal energy series $D=\sum_p(\cren)^2$
(Lemma~\ref{lem:D-conv}).
The normalised Gaussian correlation $K_\varepsilon$ localises
the diagonal first-prime contribution to $V_p(\sigma)$ by an
algebraic identity (Observation~\ref{obs:diagonal}).
Lemma~\ref{lem:ren-scale} and Observation~\ref{obs:sawtooth}
establish the renormalised scale and sawtooth structure
of~$\Hren$.
These are unconditional mathematical statements.

This paper \emph{establishes numerically}, at reference
parameters, the diagonal energy value $D\approx 9.470$
(Observation~\ref{obs:D-value}) and the finite-grid comparison
$\Zren\ge\Hlocal^{\mathrm{pr}}$ on the eight-point grid
$\kappa\in\{10,20,29,50,100,200,500,1000\}$ at every tested
point except $\kappa=20$ (Observation~\ref{obs:Z-vs-H}).
Both numerical statements are finite computations from the
stated parameters and are labelled [numerical] in their theorem
headings.

This paper \emph{leaves open} three well-posed questions named in
\S\,\ref{sec:open}:
Problem~\ref{op:weil} (Weil-normalised embedding),
Problem~\ref{op:spectral-ineq} (Spectral inequality), and
Problem~\ref{op:offdiag} (Off-diagonal control).
All three are formulated entirely within classical analytic number
theory.

% ============================================================
% §2 Construction
% ============================================================
\section{Construction of the admissible test family}
\label{sec:construct}

Recall from~\cite[\S\,2, finite-place contribution]{Paper1} that
\[
V_p(\sigma)
\;=\; 4(\log p)^2\,\frac{p^{-2\sigma}}{(1-p^{-2\sigma})^2}.
\]
Define the raw weight
\[
c_{p,\sigma}
\;:=\; V_p(\sigma)^{1/2}
\cdot \frac{p^{1/4}}{(\log p)^{1/2}},
\]
and let
\[
\varphi_\varepsilon(x)
\;=\; \frac{1}{\varepsilon\sqrt{2\pi}}\,e^{-x^2/(2\varepsilon^2)}
\]
be the standard Gaussian bump.
For fixed $\kappa$ set
\[
f_{\sigma,\varepsilon}(x)
\;:=\; \sum_{p\le\kappa} c_{p,\sigma}\,
       \varphi_\varepsilon(x-\log p).
\]
The admissible test function is obtained by antisymmetrisation:
\begin{equation}\label{eq:gstar}
g^*_{\sigma,\varepsilon}(x)
\;:=\; f_{\sigma,\varepsilon}(x)
       \;-\; f_{\sigma,\varepsilon}(-x).
\end{equation}
Under the Fourier convention
$\widehat h(t) = \int_{\mathbb R} h(x)\,e^{-itx}\,dx$,
the Gaussian Fourier pair
$\widehat\varphi_\varepsilon(t)=e^{-\varepsilon^2 t^2/2}$
gives
\begin{equation}\label{eq:FT-gstar}
\widehat g^*_{\sigma,\varepsilon}(t)
\;=\; -2i\,e^{-\varepsilon^2 t^2/2}
\sum_{p\le\kappa} c_{p,\sigma}\sin(t\log p),
\end{equation}
and hence
\begin{equation}\label{eq:gstar-modsq}
\bigl|\widehat g^*_{\sigma,\varepsilon}(\gamma_k)\bigr|^2
\;=\; 4\,e^{-\varepsilon^2\gamma_k^2}
\Bigl(\sum_{p\le\kappa}
c_{p,\sigma}\sin(\gamma_k\log p)\Bigr)^{\!2}.
\end{equation}

The admissibility class $\Sad$ used here is the one appearing in
Lagarias' formulation of Weil positivity~\cite{Lagarias}, built
on Weil's explicit formula~\cite{Weil}.
The foundational explicit-formula and test-function framework
is closely related to Barner's formulation of Weil's explicit
formula~\cite{Barner81}.
Avdispahi\'c--Smajlovi\'c extend the formula to broader
test-function classes~\cite{AvdSmaj05}; our construction remains
inside the Lagarias admissible class.
A test function $g$ is
admissible if $g$ and $\widehat g$ are both Schwartz,
$g(0)=\widehat g(0)=0$, and the prime-side and zero-side
expressions in the Weil explicit formula converge absolutely.
Explicit admissible families in the Burnol--Sonin
framework are studied in~\cite{Burnol,BurnolJTNB}.
To the best of our knowledge,
the Gaussian family $g^*_{\sigma,\varepsilon}$ used here gives a
prime-indexed realisation with closed-form weights specifically
targeting $\Hlocal$.
For fixed $\kappa$, we write $g^*_{\sigma,\varepsilon}$ for
$g^*_{\sigma,\varepsilon,\kappa}$, suppressing $\kappa$ when
no confusion arises.

\begin{lemma}[Admissibility; proved]\label{lem:admissible}
For every finite cutoff $\kappa\ge 2$, every $\sigma>0$, and
every $\varepsilon>0$, the function
$g^*_{\sigma,\varepsilon,\kappa}$ belongs to the admissible
class $\Sad$
used by Lagarias~\cite{Lagarias} in the formulation of Weil
positivity.
\end{lemma}

\begin{proof}
The conditions of Lagarias~\cite{Lagarias} --- Schwartz class
($g^*$ and $\widehat g^*$ rapidly decreasing), vanishing at
the origin ($\widehat g^*(0)=g^*(0)=0$), and absolute
convergence of the prime-side and zero-side sums --- are all
satisfied.
The Gaussian bumps $\varphi_\varepsilon$ are Schwartz, and a
finite linear combination of translated Schwartz functions is
again Schwartz; antisymmetrisation preserves the Schwartz class.
Hence $g^*_{\sigma,\varepsilon}$ and
$\widehat g^*_{\sigma,\varepsilon}$ are rapidly decreasing.
By antisymmetry, $\widehat g^*_{\sigma,\varepsilon}(0)=0$ and
$g^*_{\sigma,\varepsilon}(0)=0$.
The Gaussian decay $e^{-\varepsilon^2\gamma_k^2}$ in the
Fourier representation ensures absolute convergence of
$\sum_k|\widehat g^*_{\sigma,\varepsilon}(\gamma_k)|^2$
over all zero ordinates~$\gamma_k$.
More precisely, for a non-trivial zero $\rho=\beta+i\gamma$
with $0<\beta<1$, let
$z_\rho=(\rho-\tfrac12)/i=\gamma-i(\beta-\tfrac12)$.
Since $|\mathrm{Im}\,z_\rho|\le\tfrac12$, the Fourier
formula~\eqref{eq:FT-gstar} and the finiteness of the prime
sum give a constant $C_{\sigma,\varepsilon,\kappa}$ such that
$|\widehat g^*_{\sigma,\varepsilon}(z_\rho)|
\le C_{\sigma,\varepsilon,\kappa}\,e^{-\varepsilon^2\gamma^2/2}$.
The Riemann--von~Mangoldt bound $N(T)=O(T\log T)$ then implies
absolute convergence of the zero-side sum.
No critical-line localisation is used.
For the prime side, use the explicit Gaussian factor
in~\eqref{eq:FT-gstar}, not only the Schwartz property.
For $u=m\log p$,
\[
|\widehat g^*_{\sigma,\varepsilon}(u)|
\;\le\; C_{\sigma,\kappa}\,e^{-\varepsilon^2 u^2/2}
\;=\; C_{\sigma,\kappa}\,e^{-\varepsilon^2 m^2(\log p)^2/2}.
\]
Hence for every $B>0$ there is $C_B$ such that
$|\widehat g^*_{\sigma,\varepsilon}(m\log p)|\le C_B\,p^{-Bm}$.
Choosing $B$ large enough for the Lagarias prime-power weights
gives absolute convergence of the double prime-power sum.
\end{proof}

% ============================================================
% §3 Formal prime-side localisation
% ============================================================
\section{Formal prime-side localisation and the open Weil bridge}
\label{sec:prime-side}

Define the \emph{normalised Gaussian correlation}
\[
K_\varepsilon(x)
\;:=\;
\frac{[\varphi_\varepsilon*\check\varphi_\varepsilon](x)}
     {[\varphi_\varepsilon*\check\varphi_\varepsilon](0)}
\;=\; \exp\!\Bigl(-\frac{x^2}{4\varepsilon^2}\Bigr).
\]
Thus $K_\varepsilon(0)=1$ and $K_\varepsilon(x)\to 0$
exponentially for every fixed $x\ne 0$ as $\varepsilon\to 0$
(Lemma~\ref{lem:gauss}).

\begin{lemma}[Gaussian localisation; proved]\label{lem:gauss}
For the family $g^*_{\sigma,\varepsilon}$
of~\S\,\ref{sec:construct}, at fixed $(\sigma,\kappa)$ as
$\varepsilon\to 0$,
$K_\varepsilon(x)\to 0$ exponentially for every
fixed $x\ne 0$.
\end{lemma}

\begin{proof}
By definition,
$K_\varepsilon(x)=e^{-x^2/(4\varepsilon^2)}$,
so for fixed $x\ne 0$ and $\varepsilon\to 0$ the exponent
$-x^2/(4\varepsilon^2)\to -\infty$.
\end{proof}

\begin{observation}[Formal diagonal localisation; proved]
\label{obs:diagonal}
With the normalised Gaussian correlation $K_\varepsilon$ and the
weights $c_{p,\sigma}=V_p(\sigma)^{1/2}\,p^{1/4}/(\log p)^{1/2}$,
the diagonal first-prime contribution satisfies the algebraic identity
\[
\frac{\log p}{\sqrt{p}}\,c_{p,\sigma}^2\,K_\varepsilon(0)
\;=\; V_p(\sigma).
\]
Every fixed off-diagonal prime contribution
($p\ne q$) and every fixed higher-prime-power
contribution with non-zero logarithmic separation
satisfies $K_\varepsilon(\cdot)\to 0$ exponentially
as $\varepsilon\to 0$ at fixed~$\kappa$
(Lemma~\ref{lem:gauss}).
No uniform summation over the full prime-power side
is asserted here; that requires the separate
tail/dominated-convergence argument isolated in
Problem~\ref{op:weil}.
\end{observation}

\begin{proof}
By the definition of $c_{p,\sigma}$,
\[
\frac{\log p}{\sqrt{p}}\,c_{p,\sigma}^2
\;=\;
\frac{\log p}{\sqrt{p}}\cdot
V_p(\sigma)\,\frac{\sqrt{p}}{\log p}
\;=\; V_p(\sigma).
\]
Since $K_\varepsilon(0)=1$, the identity follows.
The off-diagonal and prime-power decay is immediate from
$K_\varepsilon(x)=e^{-x^2/(4\varepsilon^2)}\to 0$ for $x\ne 0$.
\end{proof}

Observation~\ref{obs:diagonal} shows that the construction is
structured to target $\Hlocal(\sigma,\kappa)$ on the prime side:
the weights $c_{p,\sigma}$ are precisely designed so that the
normalised diagonal contribution reproduces $V_p(\sigma)$ at
$K_\varepsilon(0)=1$, while all contributions with non-zero logarithmic
separation are exponentially suppressed.
The identity holds by construction: the content of
Observation~\ref{obs:diagonal} is the existence of a
Schwartz-class admissible family in $\Sad$ with this
localisation property.
The normalisation factor $p^{1/4}$ in the weights
$c_{p,\sigma}$ is consistent with the Haran--Connes
local $p$-adic normalisation~\cite{Haran01,Connes}.
The precise embedding of this localisation into the
Lagarias--Weil functional~\cite{Lagarias} --- including the
correct Gaussian rescaling, the zero-side contribution, and the
archimedean term --- is an open problem (Problem~\ref{op:weil}).

More explicitly, at fixed finite $\kappa$, set
\[
\delta_\kappa
\;:=\; \min_{\substack{p\ne q \\ p,q\le\kappa}}
|\log p - \log q| \;>\; 0.
\]
Then every fixed off-diagonal contribution is
$O(\exp(-\delta_\kappa^2/(4\varepsilon^2)))$
as $\varepsilon\to 0$.
For the full prime-power side of the Weil functional, where
infinitely many prime powers occur, this termwise localisation
must be combined with a separate tail/dominated-convergence
argument; that global embedding is part of
Problem~\ref{op:weil}.

\begin{remark}
The vanishing condition $\widehat g^*_{\sigma,\varepsilon}(0)=0$
(Lemma~\ref{lem:admissible}) ensures that the constant term in
the Lagarias decomposition of the Weil quadratic functional
vanishes, leaving only the prime-side, archimedean, and zero-side
contributions.
\end{remark}

% ============================================================
% §4 Renormalised weights and convergence
% ============================================================
\section{Renormalised weights and convergence}
\label{sec:renorm}

At $\sigma=\tfrac12$ define
\[
f_p \;:=\; V_p(\tfrac12) \;-\; \frac{4(\log p)^2}{p}.
\]
A direct calculation gives
\begin{equation}\label{eq:fp}
f_p \;=\; \frac{4(\log p)^2(2p-1)}{p(p-1)^2} \;>\; 0
\qquad (p\ge 2).
\end{equation}

\begin{lemma}[Positivity of the renormalised weight; proved]\label{lem:fp-pos}
For every prime $p\ge 2$, the quantity $f_p$
in~\eqref{eq:fp} is strictly positive.
\end{lemma}

\begin{proof}
The numerator $2p-1>0$ and the denominator $p(p-1)^2>0$ for
every prime $p\ge 2$.
\end{proof}

\begin{remark}[Prime-power-tail interpretation]\label{rem:fp-tail}
Using the Lambert-series identity
$V_p(\tfrac12) = 4(\log p)^2\sum_{m\ge1}m\,p^{-m}$,
one has
\[
f_p
\;=\;
4(\log p)^2\sum_{m\ge2}m\,p^{-m}.
\]
Thus $f_p$ isolates the prime-power tail of $V_p(\tfrac12)$
after removing the leading $m=1$ term
$4(\log p)^2/p$.
\end{remark}

The renormalised weight is then simply
\begin{equation}\label{eq:cren}
\cren \;:=\; f_p^{1/2} \;>\; 0,
\end{equation}
so that $(\cren)^2 = f_p$ exactly.
Numerical values for small primes:

\begin{center}
\begin{tabular}{cccc}
\toprule
$p$ & $V_p(\tfrac12)$ & $4(\log p)^2/p$ & $f_p = (\cren)^2$ \\
\midrule
2  & 3.844 & 0.961 & 2.883 \\
3  & 3.621 & 1.609 & 2.012 \\
5  & 3.238 & 2.072 & 1.166 \\
7  & 2.945 & 2.164 & 0.781 \\
11 & 2.530 & 2.091 & 0.439 \\
\bottomrule
\end{tabular}
\end{center}

\begin{lemma}[Convergence of the diagonal energy; proved]\label{lem:D-conv}
The series
\begin{equation}\label{eq:D}
D \;:=\; \sum_p (\cren)^2
\;=\; \sum_p \frac{4(\log p)^2(2p-1)}{p(p-1)^2}
\end{equation}
converges absolutely.
\end{lemma}

\begin{proof}
For large $p$, $(\cren)^2 \sim 8(\log p)^2/p^2$.
Since $\sum_n (\log n)^2/n^2 < \infty$, absolute convergence
follows by comparison with this dominating series.
\end{proof}

\begin{observation}[Diagonal energy value; numerical]\label{obs:D-value}
The partial sum over primes below $10^7$ gives
\[
\sum_{p<10^7} f_p \;=\; 9.4702318356.
\]
For $p\ge 10^7$, the closed form~\eqref{eq:fp} satisfies
$f_p < 9(\log p)^2/p^2$, so the tail is bounded by
\[
\sum_{p\ge 10^7} f_p
\;<\; \sum_{n\ge 10^7}\frac{9(\log n)^2}{n^2}
\;<\; 2.7\times 10^{-4}.
\]
Therefore $D = \sum_p f_p \in (9.4702318,\;9.4704965)$.
In particular $D\approx 9.470$ to three decimal places
(see also~\cite{Code}).
\end{observation}

\begin{remark}[Prime-zeta-tail identity]\label{rem:D-primezeta}
Summing the prime-power-tail representation of~$f_p$
(Remark~\ref{rem:fp-tail}) over all primes gives
\[
D
\;=\;
\sum_p f_p
\;=\;
4\sum_{m\ge2}m\,\sum_p(\log p)^2\,p^{-m}.
\]
The inner prime sums $\sum_p(\log p)^2 p^{-m}$ converge
absolutely for $m\ge 2$ and equal $P''(m)$, the second
derivative of the prime-zeta function
$P(s)=\sum_p p^{-s}$, since
$P''(s) = \sum_p(\log p)^2 p^{-s}$.
This identifies $D$ as the global prime-zeta-tail constant
of the renormalisation.
\end{remark}

% ============================================================
% §5 Renormalised scale and Mertens constant
% ============================================================
\section{The renormalised scale and Mertens-type prime sums}
\label{sec:scale}

\begin{lemma}[Renormalised local scale; proved]\label{lem:ren-scale}
As $\kappa\to\infty$,
\[
\Hren(\tfrac12,\kappa)
\;:=\;
\Hlocal(\tfrac12,\kappa) - 2(\log\kappa)^2
\;=\; O(\log\kappa).
\]
In particular,
$\Hren(\tfrac12,\kappa) = o((\log\kappa)^2)$.
\end{lemma}

\begin{proof}
By the Mertens estimate
\[
A(x) \;:=\; \sum_{p\le x}\frac{\log p}{p}
\;=\; \log x + O(1),
\]
Abel summation gives
\[
\sum_{p\le\kappa}\frac{(\log p)^2}{p}
\;=\; \log(\kappa)\,A(\kappa)
- \int_2^{\kappa}\frac{A(t)}{t}\,dt
\;=\; \tfrac12(\log\kappa)^2 + O(\log\kappa).
\]
Since $\sum_p f_p = D<\infty$ by Lemma~\ref{lem:D-conv},
the identity $V_p(\tfrac12)=4(\log p)^2/p + f_p$ yields
$\Hren(\tfrac12,\kappa) = O(\log\kappa)$.
\end{proof}

The leading constant $2$ in
$\Hlocal(\tfrac12,\kappa)\sim 2(\log\kappa)^2$
arises from
$4\cdot\tfrac12 = 2$.

\begin{observation}[Sawtooth structure; proved]\label{obs:sawtooth}
For real $\kappa\ge 2$, the sum
$\sum_{p\le\kappa}V_p(\tfrac12)$
is a step function with jump $V_p(\tfrac12)$ at each
prime~$p$, while $2(\log\kappa)^2$ is continuous and strictly
increasing.
Hence $\Hren(\tfrac12,\kappa)$ jumps upward at primes and
decreases between consecutive primes.
\end{observation}

\medskip
The \emph{prime-power-tail energy} at cutoff~$\kappa$ is
\begin{equation}\label{eq:Htail}
H_{\mathrm{tail}}(\kappa)
\;:=\;
\sum_{p\le\kappa} f_p
\;=\;
\sum_{p\le\kappa}
\frac{4(\log p)^2(2p-1)}{p(p-1)^2}.
\end{equation}
By Lemma~\ref{lem:D-conv}, $H_{\mathrm{tail}}(\kappa)\nearrow D$
as $\kappa\to\infty$.
The quantity $H_{\mathrm{tail}}$ is the partial sum of the
diagonal energy: it is $H_{\mathrm{tail}}$ (not~$\Hren$) that
grows by~$f_p$ at each new prime and converges to~$D$.

The same step-function mechanism as for $\pi(x)$ is present:
$\Hren$ increases by~$V_p(\tfrac12)$ at each
prime $p\le\kappa$ and decreases between consecutive primes
as $2(\log\kappa)^2$ grows.
The average drift of $\Hren$ is controlled by classical
prime-sum estimates~\cite{Mertens}:
\[
\sum_{p\le\kappa}\frac{\log p}{p} \;=\; \log\kappa + O(1)
\qquad (\kappa\to\infty).
\]
The Mertens constant
$\gamma_M\approx 0.2615$ governs the related estimate
$\sum_{p\le\kappa}1/p = \log\log\kappa + \gamma_M + o(1)$,
connecting the renormalised scale to classical prime-sum
asymptotics.

\begin{remark}
The Mertens constant
\[
\gamma_M \;=\; \gamma + \sum_p\Bigl[\log\bigl(1-\tfrac1p\bigr)+\tfrac1p\Bigr]
\;\approx\; 0.2615,
\]
where $\gamma\approx 0.5772$ is the Euler--Mascheroni constant.
The constant $\gamma_M$ governs the harmonic prime-sum asymptotic
\[
\sum_{p\le\kappa}\frac{1}{p}
\;=\; \log\log\kappa + \gamma_M + o(1).
\]
The renormalised curvature scale considered here is controlled
instead by the weighted Mertens-type sums
$\sum_{p\le\kappa}(\log p)^2/p$ and
$\sum_{p\le\kappa}(\log p)/p$;
thus $\gamma_M$ is a related prime-sum constant, but it is not
the leading constant governing
$\Hren(\tfrac12,\kappa)$.
The same Euler--Mascheroni constant $\gamma$ enters the first Li
coefficient~\cite{Li}
\[
\lambda_1 \;=\; 1 + \tfrac{\gamma}{2} - \log 2
- \tfrac{\log\pi}{2} \;\approx\; 0.0231.
\]
\end{remark}

\begin{remark}
The renormalised construction clarifies the scale of the
prime-side contribution but does not by itself resolve the
spectral inequality of Problem~\ref{op:spectral-ineq}.
\end{remark}

% ============================================================
% §6 Numerical comparison
% ============================================================
\section{Numerical comparison of zero and prime sides}
\label{sec:numerical}

We compare the renormalised zero sum with the
prime-only renormalised local curvature on a finite grid in
$\kappa$.
For the numerical comparison, we use a simplified
position-space proxy family
\[
g^{\mathrm{ren}}_{\kappa,\varepsilon}(x)
\;:=\; \sum_{p\le\kappa} \cren\,
\bigl[\varphi_\varepsilon(x-\log p)
      -\varphi_\varepsilon(x+\log p)\bigr],
\]
with weights $\cren=f_p^{1/2}$ in place of the abstract
weights~$c_{p,\sigma}$ of the family
$g^*_{\sigma,\varepsilon}$ constructed
in~\S\,\ref{sec:construct}.
This proxy family differs from $g^*_{\sigma,\varepsilon}$
in its weights; Problem~\ref{op:spectral-ineq} concerns the
original family.

In this section, the comparison is performed on the prime side
only.  Define the \emph{prime-only renormalised curvature}
\[
\Hlocal^{\mathrm{pr}}(\tfrac12,\kappa)
\;:=\; \sum_{p\le\kappa} V_p(\tfrac12) \;-\; 2(\log\kappa)^2,
\]
separating the archimedean contribution
$\psi_1(\tfrac12)\approx 2.150$, which enters the full
renormalised curvature $\Hren(\tfrac12,\kappa) =
\psi_1(\tfrac12) + \sum_{p\le\kappa}V_p(\tfrac12)
- 2(\log\kappa)^2 = \Hlocal^{\mathrm{pr}}(\tfrac12,\kappa)
+ \psi_1(\tfrac12)$
but is omitted from the grid comparison below.
The renormalised zero sum is
\[
\Zren(\kappa)
\;:=\; Z_{\mathrm{ord}}(g^{\mathrm{ren}}_{\kappa,\varepsilon};\,N),
\]
computed at $\varepsilon=0.05$, $N=100$ zeros.
The ordinate proxy
$Z_{\mathrm{ord}}(g^{\mathrm{ren}}_{\kappa,\varepsilon};\,N)$
is bounded at fixed
$\varepsilon$ and requires no subtraction of $2(\log\kappa)^2$;
that renormalisation applies to the prime
side~$\Hlocal^{\mathrm{pr}}$ only.
By the antisymmetric Gaussian construction and the Lagarias
convention~\cite{Lagarias},
\begin{equation}\label{eq:Zren-fourier}
\bigl|\widehat g^{\mathrm{ren}}_{\kappa,\varepsilon}
(\gamma_k)\bigr|^2
\;=\; 4\,e^{-\varepsilon^2\gamma_k^2}
\Bigl(\sum_{p\le\kappa}
\cren\sin(\gamma_k\log p)\Bigr)^{\!2},
\end{equation}
which is non-negative and reproduces the phasor-damping
factor $e^{-\varepsilon^2\gamma_k^2}$ of the numerical
computation.

The smoothing parameter $\varepsilon=0.05$ is required for
meaningful numerical experiments with the zero sum: for
$\varepsilon=0.5$ one has $e^{-\varepsilon^2\gamma_1^2}\approx
10^{-22}$, effectively suppressing all zero contributions, so
that no observable signal remains.
With $\cren=f_p^{1/2}$:

\begin{center}
\begin{tabular}{ccccc}
\toprule
$\kappa$ & $\Zren$ & $\Hlocal^{\mathrm{pr}}$
         & $\Zren-\Hlocal^{\mathrm{pr}}$
         & $\Zren\ge\Hlocal^{\mathrm{pr}}$ \\
\midrule
10   & 4.888 & 3.044 & $+1.844$ & $\checkmark$ \\
20   & 3.782 & 4.770 & $-0.988$ & $\times$ \\
29   & 5.520 & 3.588 & $+1.931$ & $\checkmark$ \\
50   & 5.693 & 2.881 & $+2.812$ & $\checkmark$ \\
100  & 5.311 & 1.562 & $+3.748$ & $\checkmark$ \\
200  & 5.102 & 2.355 & $+2.747$ & $\checkmark$ \\
500  & 5.156 & 1.101 & $+4.056$ & $\checkmark$ \\
1000 & 4.289 & 0.426 & $+3.863$ & $\checkmark$ \\
\bottomrule
\end{tabular}
\end{center}

\begin{observation}[Finite-grid stability of the comparison; numerical]\label{obs:Z-vs-H}
On the eight-point grid
$\kappa\in\{10,20,29,50,100,200,500,1000\}$, the inequality
$\Zren\ge\Hlocal^{\mathrm{pr}}$ holds at every tested point
except $\kappa=20$,
where $\Zren-\Hlocal^{\mathrm{pr}}=-0.988$.
The ratio $\Zren/\Hlocal^{\mathrm{pr}}$ equals about $1.5$ at
$\kappa=29$ and about $10.1$ at $\kappa=1000$.
No claim is made about $\kappa$ outside the eight tested points.
\end{observation}

The single grid violation $\Zren<\Hlocal^{\mathrm{pr}}$ at
$\kappa=20$ illustrates
the sawtooth behaviour of $\Hlocal^{\mathrm{pr}}$: directly
after a prime jump
$\Hlocal^{\mathrm{pr}}$ rises by $V_p(\tfrac12)$, then
decreases continuously until the next prime.
The pattern observed on the grid for $\kappa\ge 29$ is consistent
with the spectral inequality of Problem~\ref{op:spectral-ineq},
but does not prove it: Problem~\ref{op:spectral-ineq} concerns the unconditional
Lagarias--Weil zero-side contribution
$\mathcal{Z}_W(g^*_{1/2,\varepsilon})$, whereas the present
computation uses the finite positive-ordinate proxy
$Z_{\mathrm{ord}}$.
No positivity or inequality for $\mathcal{Z}_W$ is inferred
from the finite-grid data.

All numerical claims of this section are reproducible from the
verification scripts in the GitHub repository~\cite{Code}.

% ============================================================
% §7 Open Problems
% ============================================================
\section{Open Problems}
\label{sec:open}

The present work identifies a formal diagonal localisation
pattern (Observation~\ref{obs:diagonal}) and reduces the
question of prime-side derivation of Weil positivity for the
family $g^*_{\sigma,\varepsilon}$ to three well-posed
subproblems.
We name each and indicate the tools we expect to be relevant.

\begin{openproblem}[Weil-normalised embedding]\label{op:weil}
Determine whether the normalised diagonal localisation of
Observation~\ref{obs:diagonal} can be implemented inside the
Lagarias--Weil functional~\cite{Lagarias} for the position-space
family $g^*_{\sigma,\varepsilon}$, including:
\begin{enumerate}
\item[(a)] the correct Gaussian rescaling so that the convolution
value at the origin matches the prime-side Euler factor
$V_p(\sigma)$;
\item[(b)] the archimedean Weil-factor contribution
$\mathcal{A}_\infty(g^*_{\sigma,\varepsilon}
*\check g^*_{\sigma,\varepsilon})$,
including the identification and, if necessary, subtraction or
rescaling of the singular self-correlation term
$[\varphi_\varepsilon*\check\varphi_\varepsilon](0)
\sim (2\sqrt\pi\,\varepsilon)^{-1}$
as $\varepsilon\to 0$;
\item[(c)] the zero-side contribution
$\mathcal{Z}_W(g^*_{\sigma,\varepsilon})$ in its unconditional
form, without assuming critical-line localisation.
\end{enumerate}
A positive answer would establish the formal prime-side identity
\[
W(g^*_{\sigma,\varepsilon}*\check g^*_{\sigma,\varepsilon})
\;=\; \mathcal{Z}_W(g^*_{\sigma,\varepsilon})
\;-\;\Hlocal(\sigma,\kappa)
\;+\; R(\sigma,\varepsilon,\kappa)
\]
as a proved statement, with $R\to 0$ as $\varepsilon\to 0$
at fixed $(\sigma,\kappa)$.

For the global (ad\`ele) setting, Connes and
Consani~\cite{ConnesConsani21} establish Weil positivity at
the archimedean place via a trace formula on Sonin's space,
using prolate spheroidal wave functions and Toeplitz matrix
theory; the UV prolate spectrum is further developed
in~\cite{ConnesMoscovici22}.
Our finite-cutoff Gaussian family
$g^*_{\sigma,\varepsilon}\in\Sad$ operates in the
Lagarias local framework at fixed~$\kappa$; the archimedean
limit as $\varepsilon\to 0$ in this setting remains the
content of~(b).
Thus \cite{ConnesConsani21} provides an archimedean positivity
precedent, but not a direct evaluation of
$\mathcal{A}_\infty(g^*_{\sigma,\varepsilon}
*\check g^*_{\sigma,\varepsilon})$
for the finite-cutoff Gaussian family used here; in particular,
it does not control the $\varepsilon\to 0$ singular
self-correlation term in the Lagarias normalisation.
\end{openproblem}

\begin{openproblem}[Spectral inequality on the critical line]\label{op:spectral-ineq}
Prove that for every cutoff $\kappa$ there exists
$\varepsilon_0(\kappa)>0$ such that for all
$0<\varepsilon<\varepsilon_0(\kappa)$,
\begin{equation}\label{eq:openineq}
\mathcal{Z}_W(g^*_{1/2,\varepsilon})
\;\ge\; \Hlocal(\tfrac12,\kappa),
\end{equation}
where $\mathcal{Z}_W$ denotes the unconditional zero-side
contribution in the Lagarias--Weil functional
(Problem~\ref{op:weil}(c)).
This inequality would be a necessary component of
a proof of Weil positivity through this specific family,
after the embedding problem of Problem~\ref{op:weil} has been
resolved.
By itself it is not an RH criterion and is not
asserted to be necessary for RH in general.
By Lagarias~\cite{Lagarias}, the equivalence
$\mathrm{RH}\iff W(f*\check f)\ge 0$ holds for all
$f\in\Sad$.
We note that Bombieri~\cite{Bombieri2000} analyses Weil's
quadratic functional on support-restricted $L^2$-spaces and
finite-dimensional truncations.
This illustrates that positivity is a delicate statement about
the chosen test-function class and limiting process: on the full
relevant admissible class, Weil positivity is RH-equivalent,
not a formal consequence of the local sign pattern or of the
diagonal localisation in Observation~\ref{obs:diagonal}.
The expected tools are: explicit lower bounds on the zero-side
quadratic form by the unconditional zero-counting density;
finite-$\kappa$ cross-term estimates for the original family
$g^*_{1/2,\varepsilon}$.
Problem~\ref{op:offdiag} formulates a related but distinct
proxy-family question in the joint regime
$\kappa\to\infty$, $\varepsilon=\varepsilon(\kappa)\to 0$.
On the tested grid the inequality
$\Zren\ge\Hlocal^{\mathrm{pr}}$ holds at every
tested point except $\kappa=20$ (Observation~\ref{obs:Z-vs-H}),
consistent with~\eqref{eq:openineq} but not establishing it
analytically.
\end{openproblem}

\begin{openproblem}[Off-diagonal control of the zero sum]\label{op:offdiag}
To avoid conflict with the finite numerical proxy
$Z_{\mathrm{ord}}(g;N)$ used in~\S\,\ref{sec:numerical},
we write $Z_+^{\infty}(g)$ for the infinite positive-ordinate
quadratic form obtained by summing over all $\gamma_k>0$.
Thus $Z_{\mathrm{ord}}(g;N)$ is finite and numerical,
whereas $Z_+^{\infty}(g)$ is the asymptotic object in this
problem.

Prove unconditionally that the combined non-diagonal
ratio contribution and product-term remainder in the
positive-ordinate zero sum are asymptotically negligible
compared with the diagonal $R_1$ contribution.
Set
\[
I_+(\varepsilon)
\;:=\; \sum_{\gamma_k>0} e^{-\varepsilon^2\gamma_k^2}.
\]
By the Riemann--von~Mangoldt zero-counting formula and partial
summation,
\[
I_+(\varepsilon)
\;=\;
\frac{1}{4\sqrt\pi\,\varepsilon}\,\log\frac{1}{\varepsilon}
\;+\; O\!\left(\frac{1}{\varepsilon}\right)
\qquad (\varepsilon\to 0^+).
\]
Thus, for
$\varepsilon(\kappa)\sim (4\sqrt\pi\,\log\kappa)^{-1}$,
one has
$I_+(\varepsilon(\kappa))\sim\log\kappa\,\log\log\kappa$.
For
\[
R_a(\varepsilon)
\;:=\; \sum_{\gamma_k>0}
e^{-\varepsilon^2\gamma_k^2}\cos(\gamma_k\log a),
\]
one has $R_1(\varepsilon)=I_+(\varepsilon)$.
The identity
\[
Z_+^{\infty}(g^{\mathrm{ren}}_{\kappa,\varepsilon})
\;=\; 2\!\sum_{p,q\le\kappa}
\cren\,c_q^{\mathrm{ren}}\,R_{p/q}(\varepsilon)
\;-\; 2\!\sum_{p,q\le\kappa}
\cren\,c_q^{\mathrm{ren}}\,R_{pq}(\varepsilon)
\]
follows from
$4\sin(\gamma\log p)\sin(\gamma\log q)
= 2\cos(\gamma\log(p/q)) - 2\cos(\gamma\log(pq))$.
The principal diagonal term is
\[
2\sum_{p\le\kappa}(\cren)^2\,I_+(\varepsilon)
\;=\; 2\,H_{\mathrm{tail}}(\kappa)\,I_+(\varepsilon).
\]
Prove unconditionally that, in the expansion of
$Z_+^{\infty}(g^{\mathrm{ren}}_{\kappa,\varepsilon})$
above, for the scale
\[
\varepsilon(\kappa)
\;\sim\; (4\sqrt\pi\,\log\kappa)^{-1},
\]
the combined non-diagonal and product-term remainder
satisfies
\[
2\!\!\!\sum_{\substack{p,q\le\kappa \\ p\ne q}}\!\!\!
\cren\,c_q^{\mathrm{ren}}\,R_{p/q}(\varepsilon)
\;-\; 2\!\sum_{p,q\le\kappa}
\cren\,c_q^{\mathrm{ren}}\,R_{pq}(\varepsilon)
\;=\; o\bigl(H_{\mathrm{tail}}(\kappa)\,I_+(\varepsilon)\bigr).
\]
Then, since $H_{\mathrm{tail}}(\kappa)\to D$
(Lemma~\ref{lem:D-conv}), one obtains
\[
Z_+^{\infty}(g^{\mathrm{ren}}_{\kappa,\varepsilon(\kappa)})
\;\sim\; 2D\,I_+(\varepsilon(\kappa)).
\]
For example, if
$\varepsilon(\kappa)\sim (4\sqrt\pi\,\log\kappa)^{-1}$,
then $I_+(\varepsilon(\kappa))\sim\log\kappa\,\log\log\kappa$
and hence
\[
Z_+^{\infty}(g^{\mathrm{ren}}_{\kappa,\varepsilon(\kappa)})
\;\sim\; 2D\log\kappa\log\log\kappa,
\]
motivated by Selberg-type orthogonality and mean-square
considerations~\cite{Selberg}.
The expected tools are: analytic continuation of
$R_a(\varepsilon)$ for $a\ne 1$ via classical
explicit-formula techniques; cancellation by oscillation, with
the $R_{pq}$ product terms providing additional cancellation.
The Beurling--Selberg extremal function framework used by
Carneiro--Chandee--Littmann--Milinovich~\cite{CCLM17}
and the semidefinite programming methods of
Chirre--Gon\c{c}alves--de~Laat~\cite{CGdL20}
provide methodological precedents for bounding
zero-statistical sums with Gaussian test functions.
They are cited here as analytic tools to be adapted,
not as hypotheses or as unconditional inputs to
Problem~\ref{op:offdiag}.
This problem is formulated entirely within classical analytic
number theory.
\end{openproblem}

\begin{remark}
Problems~\ref{op:weil}, \ref{op:spectral-ineq},
and~\ref{op:offdiag} address
different regimes.
Problem~\ref{op:weil} concerns the structural embedding of the
diagonal localisation into the Lagarias--Weil functional.
Problem~\ref{op:spectral-ineq} is a fixed-cutoff, small-$\varepsilon$
inequality: for each finite~$\kappa$, it asks whether there exists
$\varepsilon_0(\kappa)>0$ such that the inequality holds for all
$0<\varepsilon<\varepsilon_0(\kappa)$.
Problem~\ref{op:offdiag} concerns the joint limit
$\kappa\to\infty$, $\varepsilon=\varepsilon(\kappa)\to 0$, and
addresses the asymptotic structure of the ordinate sum in that
regime.
The diagonal scale $2D\log\kappa\log\log\kappa$ in
Problem~\ref{op:offdiag} is asymptotically smaller than
$\Hlocal(\tfrac12,\kappa)\sim 2(\log\kappa)^2$; the two
problems are therefore complementary rather than directly
linked.
\end{remark}

% ============================================================
% §8 Non-Circularity
% ============================================================
\section*{Non-Circularity}
\label{sec:nocirc}

We document explicitly that the results of this paper do not
rely on the objects they are directed toward.

\textbf{No RH as input.}
At no point is the Riemann Hypothesis assumed.
The non-trivial zeros $\rho$ enter the construction only through
the classical zero-counting density $N(T)\sim T\log T/(2\pi)$,
which is unconditional, and through finite Gaussian-smoothed
sums for which the symmetric pairing $\rho\leftrightarrow\bar\rho$
is used without further hypothesis.
The reference point $\sigma=\tfrac12$ is chosen as the object
of study, not assumed to be the unique zero location.

\textbf{No GUE, Montgomery, or Hilbert--P\'olya hypothesis.}
The construction of $g^*_{\sigma,\varepsilon}$ uses only
Schwartz-class smoothing and antisymmetrisation; it makes no
reference to random-matrix statistics, pair-correlation
conjectures, or any postulated spectral realisation of the
zeros.
The diagonal scale in Problem~\ref{op:offdiag} is obtained by
combining $D=\sum_p f_p$ with the unconditional
Riemann--von~Mangoldt zero-counting density; no GUE,
Montgomery pair-correlation, or Hilbert--P\'olya input is used.

\textbf{The Weil explicit formula is classical.}
The diagonal localisation (Observation~\ref{obs:diagonal})
is an algebraic identity in the weights $c_{p,\sigma}$,
independent of the Weil functional.
The embedding into the Lagarias--Weil
framework~\cite{Lagarias} is formulated as an open problem
(Problem~\ref{op:weil}); no unverified identity is used.

\textbf{Numerical results are clearly marked.}
The diagonal energy value $D\approx 9.470$ and the
comparison-grid statements of~\S\,\ref{sec:numerical} rely on
direct computation; each statement is labelled [numerical] in
its heading.

\textbf{Open problems are named, not used as hypotheses.}
Problems~\ref{op:spectral-ineq} and~\ref{op:offdiag} are stated
openly in~\S\,\ref{sec:open}.
No claim of this paper is conditional on a positive answer to
either problem.

% ============================================================
% Acknowledgments
% ============================================================
\section*{Acknowledgments}

The author thanks Jean-Fran\c{c}ois Burnol for confirming that
the curvature decomposition of~\cite{Paper1} does not appear in
the literature known to him, and the readers who provided
feedback on earlier versions of this paper.

% ============================================================
% Call for Feedback
% ============================================================
\section*{Call for Feedback}

This paper is part of an ongoing open research initiative
toward the Riemann Hypothesis.
All papers in the series are freely available on Zenodo under
CC~BY~4.0, and the accompanying verification code is hosted on
GitHub.

The author welcomes critical feedback, corrections, and
collaboration:
{\sloppy
Zenodo (\url{https://zenodo.org/search?q=Ulrich+Tehrani}) and
GitHub (\url{https://github.com/utehrani/analysislab-nt}).
\par}

Topics of particular interest include: analytical approaches to
the spectral inequality of Problem~\ref{op:spectral-ineq};
unconditional control of the off-diagonal cross-terms of
Problem~\ref{op:offdiag}; and the relation between the diagonal
constant $D$ in~\eqref{eq:D} and the classical zero-counting
density entering $I_+(\varepsilon)$.

We particularly welcome expert comment on the following
question: for the Gaussian-smoothed proxy family
$g^{\mathrm{ren}}_{\kappa,\varepsilon}$ used
in~\S\,\ref{sec:numerical},
is the asymptotic comparison
\[
Z_+^{\infty}(g^{\mathrm{ren}}_{\kappa,\varepsilon(\kappa)})
\;\sim\; 2D\log\kappa\log\log\kappa,
\qquad
\varepsilon(\kappa)\sim (4\sqrt\pi\,\log\kappa)^{-1},
\]
with $D$ as in~\eqref{eq:D}, accessible by classical sieve or
explicit-formula techniques without recourse to random-matrix
statistics?

% ============================================================
% References
% ============================================================
\newpage
\begin{thebibliography}{99}

\bibitem{Paper1}
U.~Tehrani,
``A Curvature Decomposition of the Explicit Formula
for the Riemann Zeta Function'',
Zenodo, \url{https://doi.org/10.5281/zenodo.19025598}, 2026.

\bibitem{Code}
U.~Tehrani,
\emph{Verification scripts for the prime--zero coupling series},
GitHub, \url{https://github.com/utehrani/analysislab-nt}.

\bibitem{Lagarias}
J.C.~Lagarias,
``On a positivity property of the Riemann $\xi$-function'',
\emph{Acta Arithmetica} \textbf{89} (1999), 217--234.

\bibitem{Weil}
A.~Weil,
``Sur les `formules explicites' de la th\'eorie des nombres
premiers'',
\emph{Comm.\ S\'em.\ Math.\ Univ.\ Lund},
tome suppl\'ementaire d\'edi\'e \`a M.~Riesz, 1952,
252--265.

\bibitem{Li}
X.-J.~Li,
``The positivity of a sequence of numbers and the Riemann
hypothesis'',
\emph{J.\ Number Theory} \textbf{65}.2 (1997), 325--333;
doi:10.1006/jnth.1997.2137.

\bibitem{Mertens}
F.~Mertens,
``Ein Beitrag zur analytischen Zahlentheorie'',
\emph{J.\ Reine Angew.\ Math.} \textbf{78} (1874), 46--62.

\bibitem{Selberg}
A.~Selberg,
``Contributions to the theory of the Riemann zeta-function'',
\emph{Arch.\ Math.\ Naturvid.} \textbf{48}.5 (1946), 89--155.

\bibitem{Connes}
A.~Connes,
``Trace Formula in Noncommutative Geometry and the Zeros of the
Riemann Zeta Function'',
\emph{Selecta Math.\ (N.S.)} \textbf{5}.1 (1999), 29--106.

\bibitem{BerryKeating}
M.V.~Berry and J.P.~Keating,
``The Riemann Zeros and Eigenvalue Asymptotics'',
\emph{SIAM Review} \textbf{41}.2 (1999), 236--266.

\bibitem{deBranges}
L.~de~Branges,
``The Riemann hypothesis for Hilbert spaces of entire functions'',
\emph{Bull.\ Amer.\ Math.\ Soc.\ (N.S.)} \textbf{15}.1 (1986),
1--17.

\bibitem{Burnol}
J.-F.~Burnol,
``On Fourier and Zeta(s)'',
\emph{Forum Math.} \textbf{16}.6 (2004), 789--840.

\bibitem{ConnesConsani21}
A.~Connes and C.~Consani,
``Weil positivity and Trace formula the archimedean place'',
\emph{Selecta Math.\ (N.S.)} \textbf{27}.4 (2021), article~77;
doi:10.1007/s00029-021-00689-4.

\bibitem{ConnesMoscovici22}
A.~Connes and H.~Moscovici,
``The UV prolate spectrum matches the zeros of zeta'',
\emph{Proc.\ Natl.\ Acad.\ Sci.\ USA} \textbf{119}.22 (2022),
e2123174119.

\bibitem{Barner81}
K.~Barner,
``On A.~Weil's explicit formula'',
\emph{J.\ Reine Angew.\ Math.} \textbf{323} (1981), 139--152.

\bibitem{Bombieri2000}
E.~Bombieri,
``Remarks on Weil's quadratic functional in the theory
of prime numbers, I'',
\emph{Rend.\ Mat.\ Acc.\ Lincei}, Ser.~IX,
\textbf{11} (2000), 183--233.

\bibitem{BurnolJTNB}
J.-F.~Burnol,
``An adelic causality problem related to abelian
$L$-functions'',
\emph{J.\ Th\'eor.\ Nombres Bordeaux} \textbf{16}.1
(2004), 65--94.

\bibitem{Haran01}
S.~Haran,
\emph{The Mysteries of the Real Prime},
LMS Monographs \textbf{25}, Oxford University Press, 2001.

\bibitem{AvdSmaj05}
M.~Avdispahi\'c and L.~Smajlovi\'c,
``Explicit formula for a fundamental class of functions'',
\emph{Bull.\ Belg.\ Math.\ Soc.\ Simon Stevin}
\textbf{12}.4 (2005), 569--587.

\bibitem{CCLM17}
E.~Carneiro, V.~Chandee, F.~Littmann,
and M.~B.~Milinovich,
``Hilbert spaces and the pair correlation of zeros
of the Riemann zeta-function'',
\emph{J.\ Reine Angew.\ Math.\ (Crelle)}
\textbf{725} (2017), 143--182.

\bibitem{CGdL20}
A.~Chirre, F.~Gon\c{c}alves, and D.~de~Laat,
``Pair correlation estimates for the zeros of the
zeta function via semidefinite programming'',
\emph{Adv.\ Math.} \textbf{361} (2020),
article~106926.

\end{thebibliography}

% ============================================================
% Footer — normative format
% ============================================================
\enlargethispage{3\baselineskip}
\vspace*{\fill}
\begin{minipage}{\linewidth}
\rule{\linewidth}{0.4pt}\smallskip\\
\small\emph{Paper~2 v36 \textperiodcentered\ May~2026
\textperiodcentered\ Pauca sed matura}
\end{minipage}

\end{document}
